% REMEMBER TO SET LANGUAGE!
\documentclass[a4paper,10pt,english]{article}
\usepackage[utf8]{inputenc}
\usepackage[english]{babel}
% Standard stuff
\usepackage{amsmath,graphicx,babel,varioref,verbatim,amsfonts}
% colors in text
\usepackage[usenames,dvipsnames,svgnames,table]{xcolor}
% Hyper refs
\usepackage[framestyle=none,framefit=yes,heightadjust=all,framearound=all]{floatrow}    
\floatsetup[figure]{style=Boxed,framearound=all}
\usepackage{fancyhdr}
\usepackage[colorlinks]{hyperref}

% Document formatting
\setlength{\parindent}{0mm}
\setlength{\parskip}{1.5mm}
%\setcounter{section}{-1} Hvis du vil ha 0 som første avsnitt
\numberwithin{figure}{section} 
\numberwithin{table}{section}
\numberwithin{equation}{section}

%Color scheme for listings
\usepackage{textcomp}
\definecolor{listinggray}{gray}{0.9}
\definecolor{lbcolor}{rgb}{0.9,0.9,0.9}

%Listings configuration
\usepackage{listings}
\lstset{
	backgroundcolor=\color{lbcolor},
	tabsize=4,
	rulecolor=,
	language=python,
        basicstyle=\scriptsize,
        upquote=true,
        aboveskip={1.5\baselineskip},
        columns=fixed,
	numbers=left,
        showstringspaces=false,
        extendedchars=true,
        breaklines=true,
        prebreak = \raisebox{0ex}[0ex][0ex]{\ensuremath{\hookleftarrow}},
        frame=single,
        showtabs=false,
        showspaces=false,
        showstringspaces=false,
        identifierstyle=\ttfamily,
        keywordstyle=\color[rgb]{0,0,1},
        commentstyle=\color[rgb]{0.133,0.545,0.133},
        stringstyle=\color[rgb]{0.627,0.126,0.941}
        }
        
%User settings
\newcounter{subproject}
\renewcommand{\thesubproject}{\alph{subproject}}
\newenvironment{subproj}{
\begin{description}
\item[\refstepcounter{subproject}(\thesubproject)]
}{\end{description}}
\newcommand{\HRule}{\rule{\linewidth}{0.5mm}}
\newcommand{\eqs}{\begin{equation}}
\newcommand{\eqf}{\end{equation}}


%Other
\usepackage{geometry}
\usepackage{booktabs}
\newcommand{\shaderow}{\rowcolor{gray!20}[2pt][2pt]}
\usepackage[boxed,linesnumbered,lined]{algorithm2e}

%Lettering instead of numbering in different layers
%\renewcommand{\labelenumi}{\alph{enumi}}
%\renewcommand{\thesubsection}{\alph{subsection}}



\begin{document}
%Header
\pagestyle{fancy}
\renewcommand{\sectionmark}[1]{\markright{#1}{}}

\pagestyle{fancy}
\renewcommand{\sectionmark}[1]{\markright{\thesection\ #1}}

\fancyhf{}
%Header
%\rhead{\fancyplain{}{ NAVN PÅ PROSJEKT}} % predefined ()
\lhead{\fancyplain{}{\rightmark }} % 1. sectionname, 1.1 subsection name etc
\cfoot{\fancyplain{}{\thepage}}


%Forside
\begin{titlepage}
\begin{center}

% Upper part of the page. The '~' is needed because \\
% only works if a paragraph has started.
\includegraphics[width=0.4\textwidth]{forside.jpg}~\\[1cm]

\textsc{\LARGE Project 3}\\[1.5cm]

\textsc{\Large FYS3150 - Computational physics}\\[0.5cm]

% Title
\HRule \\[0.4cm]
{ \huge \bfseries Quantum dots \\[0.4cm] }

\HRule \\[1.5cm]

% Author and supervisor
\begin{minipage}{0.4\textwidth}
\begin{flushleft} \large
\emph{Author:}\\
Vidar \textsc{Skogvoll}
\end{flushleft}
\end{minipage}
\begin{minipage}{0.4\textwidth}
\begin{flushright} \large
\emph{ } \\
 \textsc{ }
\end{flushright}
\end{minipage}

\vfill

% Bottom of the page
{\large \today}

\end{center}
\end{titlepage}
\setcounter{page}{2}

\begin{abstract}
Here is a short summary of the project.
\end{abstract}

\hypersetup{linkcolor=black}
\tableofcontents 
\hypersetup{linkcolor=red}
\clearpage




\section{Introduction}
Quantum mechanics is an exciting field. 








\section{Theory}
Here is all the theory needed to understand the project.


\subsection{The numerical foundation}
This is the section explaining the numerical theory upon which the
the project is built. 

\subsubsection{Monte Carlo methods}

\subsubsection{Importance sampling}

\subsubsection{Metropolis algorithm}


\subsection{The physical system}
This is the section explaining the physics of the system. 

\subsubsection{Quantum mechanics}

\subsubsection{Discretization}

\subsubsection{The virial theorem}








\section{Method}
This is the section explaining what has been done.


\subsection{Subsection}
These will become apparent as the work boils down. 








\section{Results and discussion}
Section listing results and discussing them. 


\subsection{Subsection}
The same as the method section. 








\section{Conclusion}

\end{document}